
\documentclass{article}
\usepackage{colortbl}
\usepackage{makecell}
\usepackage{multirow}
\usepackage{supertabular}

\begin{document}

\newcounter{utterance}

\centering \large Interaction Transcript for game `clean\_up', experiment `0\_easy\_3obj\_hu', episode 2 with Qwen3.6{-}35B{-}A3B{-}FP8{-}without{-}reasoning.
\vspace{24pt}

{ \footnotesize  \setcounter{utterance}{1}
\setlength{\tabcolsep}{0pt}
\begin{supertabular}{c@{$\;$}|p{.15\linewidth}@{}p{.15\linewidth}p{.15\linewidth}p{.15\linewidth}p{.15\linewidth}p{.15\linewidth}}
    \# & $\;$A & \multicolumn{4}{c}{Game Master} & $\;\:$B\\
    \hline

    \theutterance \stepcounter{utterance}  
    & & \multicolumn{4}{p{0.6\linewidth}}{
        \cellcolor[rgb]{0.9,0.9,0.9}{
            \makecell[{{p{\linewidth}}}]{
                \texttt{\tiny{[P1$\langle$GM]}}
                \texttt{Én vagyok a játékmestered, és egy együttműködésen alapuló játékot játszol az alábbi rács segítségével:\textbackslash{}n\textbackslash{}n```\textbackslash{}n 1234567} \\
\texttt{╔═══════╗} \\
\texttt{║◌◌◌◌L◌◌║ 1} \\
\texttt{║P◌◌◌◌◌◌║ 2} \\
\texttt{║◌C◌◌◌┌─╢ 3} \\
\texttt{╟──┐◌◌│◌║ 4} \\
\texttt{║◌◌├──┼─╢ 5} \\
\texttt{║◌◌│◌◌│◌║ 6} \\
\texttt{║◌◌│◌◌│◌║ 7} \\
\texttt{╚══╧══╧═╝} \\
\texttt{\textbackslash{}n```\textbackslash{}n\textbackslash{}n* A felső szélén az x{-}koordináták jobbra növekednek, a jobb szélén pedig az y{-}koordináták lefelé növekednek.\textbackslash{}n* A következő tárgyak véletlenszerűen helyezkednek el a rácsodon: 'C', 'L', 'P'.\textbackslash{}n\textbackslash{}n} \\
\texttt{A másik játékos a táblának egy másik változatát látja, ahol a tárgyak különböző véletlenszerű helyeken vannak. Te nem látod az ő tábláját, és ő sem látja a tiédet.} \\
\\ 
\\ 
\texttt{**Cél:**} \\
\\ 
\texttt{Mindkét játékosnak el kell helyeznie a saját tábláján lévő tárgyakat úgy, hogy azonos tárgyak ugyanazokra a koordinátákra kerüljenek. Kommunikálnotok kell, hogy közös célt érjetek el.} \\
\\ 
\\ 
\texttt{**Szabályok:**} \\
\\ 
\texttt{* Minden körben pontosan egyet küldhetsz az alábbi két parancs közül:} \\
\texttt{1. `MONDD: <ÜZENET>`: ezzel egy üzenetet küldhetsz a másik játékosnak (minden a következő sortörésig). Én továbbítom neki.} \\
\texttt{2. `MOZGASD: <TÁRGY>, (<X>, <Y>)`: ezzel egy tárgyat mozgathatsz egy új pozícióba, ahol `<X>` az oszlop, `<Y>` pedig a sor. Én visszajelzek, hogy sikeres volt{-}e a lépés.} \\
\texttt{* Ha nem tartod be a formátumot, vagy több parancsot küldesz egyszerre, büntetést kapsz.} \\
\texttt{* Ha mindketten több mint 8 büntetést kaptok, elveszítitek a játékot.} \\
\texttt{* Elengedhetetlen, hogy egyeztessetek a célállapotról! Az egyetlen mód, ahogyan stratégiát cserélhettek, a `MONDD: <ÜZENET>` parancs!} \\
\\ 
\\ 
\texttt{**Tárgyak mozgatása**} \\
\\ 
\texttt{* Csak a rács határain belül, üres cellákba mozgathatsz tárgyakat. Az üres cellák '◌' szimbólummal vannak jelölve.} \\
\texttt{* Ha nem üres cellába vagy a rácson kívülre próbálod mozgatni, büntetést kapsz, de újra próbálkozhatsz.} \\
\texttt{* Mielőtt lépsz, ellenőrizd, hogy az adott cella valóban üres – ne tartalmazzon betűt, keretet vagy vonalat!} \\
\\ 
\\ 
\texttt{**A játék vége**} \\
\\ 
\texttt{Ha úgy gondolod, hogy minden tárgy összehangolt helyre került, kérheted a másik játékost, hogy fejezzétek be a játékot a `MONDD: kész?` parancs elküldésével. Ha a másik játékos kér erre, és te `MONDD: kész!`{-}vel válaszolsz, a játék véget ér.} \\
\\ 
\texttt{A játékot megnyeritek, ha az 12 körön belül véget ér. Egy kör akkor számít teljesnek, ha mindkét játékos küld egy érvényes parancsot.} \\
\\ 
\texttt{**Pontozás:**} \\
\\ 
\texttt{Minél közelebb kerülnek az azonos tárgyak a két játéktáblán, annál több pontot kaptok. A hibákért levonás jár. Vajon meg tudjátok dönteni a rekordot?} \\
\\ 
\\ 
\texttt{Kérlek, küldj egy üzenetet a másik játékosnak a játék elindításához!} \\
            }
        }
    }
    & \\ \\

    \theutterance \stepcounter{utterance}  
    & \multicolumn{4}{p{0.6\linewidth}}{
        \cellcolor[rgb]{0.9,0.9,0.9}{
            \makecell[{{p{\linewidth}}}]{
                \texttt{\tiny{[P1$\rangle$GM]}}
                \texttt{MONDD: Szia! Kezdjük el a koordinációs játékot. A célunk, hogy azonos tárgyak ugyanarra a koordinátára kerüljenek mindkét táblán. Kérlek, jelezd, hogy melyik tárgyat (C, L vagy P) szeretnéd először mozgatni, és hová, vagy egyeztessünk egy közös stratégiát!} \\
            }
        }
    }
    & & \\ \\

    \theutterance \stepcounter{utterance}  
    & & & \multicolumn{2}{p{0.3\linewidth}}{
        \cellcolor[rgb]{0.9,0.9,0.9}{
            \makecell[{{p{\linewidth}}}]{
                \texttt{\tiny{[GM$|$GM]}}
                \texttt{MONDD: Szia! Kezdjük el a koordinációs játékot. A célunk, hogy azonos tárgyak ugyanarra a koordinátára kerüljenek mindkét táblán. Kérlek, jelezd, hogy melyik tárgyat (C, L vagy P) szeretnéd először mozgatni, és hová, vagy egyeztessünk egy közös stratégiát!} \\
            }
        }
    }
    & & \\ \\

    \theutterance \stepcounter{utterance}  
    & & \multicolumn{4}{p{0.6\linewidth}}{
        \cellcolor[rgb]{0.9,0.9,0.9}{
            \makecell[{{p{\linewidth}}}]{
                \texttt{\tiny{[GM$\rangle$P2]}}
                \texttt{Én vagyok a játékmestered, és egy együttműködésen alapuló játékot játszol az alábbi rács segítségével:\textbackslash{}n\textbackslash{}n```\textbackslash{}n 1234567} \\
\texttt{╔═══════╗} \\
\texttt{║◌◌◌◌C◌◌║ 1} \\
\texttt{║L◌◌◌◌◌◌║ 2} \\
\texttt{║P◌◌◌◌┌─╢ 3} \\
\texttt{╟──┐◌◌│◌║ 4} \\
\texttt{║◌◌├──┼─╢ 5} \\
\texttt{║◌◌│◌◌│◌║ 6} \\
\texttt{║◌◌│◌◌│◌║ 7} \\
\texttt{╚══╧══╧═╝} \\
\texttt{\textbackslash{}n```\textbackslash{}n\textbackslash{}n* A felső szélén az x{-}koordináták jobbra növekednek, a jobb szélén pedig az y{-}koordináták lefelé növekednek.\textbackslash{}n* A következő tárgyak véletlenszerűen helyezkednek el a rácsodon: 'C', 'L', 'P'.\textbackslash{}n\textbackslash{}n} \\
\texttt{A másik játékos a táblának egy másik változatát látja, ahol a tárgyak különböző véletlenszerű helyeken vannak. Te nem látod az ő tábláját, és ő sem látja a tiédet.} \\
\\ 
\\ 
\texttt{**Cél:**} \\
\\ 
\texttt{Mindkét játékosnak el kell helyeznie a saját tábláján lévő tárgyakat úgy, hogy azonos tárgyak ugyanazokra a koordinátákra kerüljenek. Kommunikálnotok kell, hogy közös célt érjetek el.} \\
\\ 
\\ 
\texttt{**Szabályok:**} \\
\\ 
\texttt{* Minden körben pontosan egyet küldhetsz az alábbi két parancs közül:} \\
\texttt{1. `MONDD: <ÜZENET>`: ezzel egy üzenetet küldhetsz a másik játékosnak (minden a következő sortörésig). Én továbbítom neki.} \\
\texttt{2. `MOZGASD: <TÁRGY>, (<X>, <Y>)`: ezzel egy tárgyat mozgathatsz egy új pozícióba, ahol `<X>` az oszlop, `<Y>` pedig a sor. Én visszajelzek, hogy sikeres volt{-}e a lépés.} \\
\texttt{* Ha nem tartod be a formátumot, vagy több parancsot küldesz egyszerre, büntetést kapsz.} \\
\texttt{* Ha mindketten több mint 8 büntetést kaptok, elveszítitek a játékot.} \\
\texttt{* Elengedhetetlen, hogy egyeztessetek a célállapotról! Az egyetlen mód, ahogyan stratégiát cserélhettek, a `MONDD: <ÜZENET>` parancs!} \\
\\ 
\\ 
\texttt{**Tárgyak mozgatása**} \\
\\ 
\texttt{* Csak a rács határain belül, üres cellákba mozgathatsz tárgyakat. Az üres cellák '◌' szimbólummal vannak jelölve.} \\
\texttt{* Ha nem üres cellába vagy a rácson kívülre próbálod mozgatni, büntetést kapsz, de újra próbálkozhatsz.} \\
\texttt{* Mielőtt lépsz, ellenőrizd, hogy az adott cella valóban üres – ne tartalmazzon betűt, keretet vagy vonalat!} \\
\\ 
\\ 
\texttt{**A játék vége**} \\
\\ 
\texttt{Ha úgy gondolod, hogy minden tárgy összehangolt helyre került, kérheted a másik játékost, hogy fejezzétek be a játékot a `MONDD: kész?` parancs elküldésével. Ha a másik játékos kér erre, és te `MONDD: kész!`{-}vel válaszolsz, a játék véget ér.} \\
\\ 
\texttt{A játékot megnyeritek, ha az 12 körön belül véget ér. Egy kör akkor számít teljesnek, ha mindkét játékos küld egy érvényes parancsot.} \\
\\ 
\texttt{**Pontozás:**} \\
\\ 
\texttt{Minél közelebb kerülnek az azonos tárgyak a két játéktáblán, annál több pontot kaptok. A hibákért levonás jár. Vajon meg tudjátok dönteni a rekordot?} \\
\\ 
\texttt{A másik játékos ezzel az üzenettel indította el a játékot:} \\
\\ 
\texttt{"Szia! Kezdjük el a koordinációs játékot. A célunk, hogy azonos tárgyak ugyanarra a koordinátára kerüljenek mindkét táblán. Kérlek, jelezd, hogy melyik tárgyat (C, L vagy P) szeretnéd először mozgatni, és hová, vagy egyeztessünk egy közös stratégiát!"} \\
\\ 
\texttt{Mi az első parancsod?} \\
            }
        }
    }
    & \\ \\

    \theutterance \stepcounter{utterance}  
    & & & \multicolumn{4}{p{0.6\linewidth}}{
        \cellcolor[rgb]{0.9,0.9,0.9}{
            \makecell[{{p{\linewidth}}}]{
                \texttt{\tiny{[GM$\langle$P2]}}
                \texttt{Mivel a célunk, hogy a tárgyak (C, L, P) azonos koordinátákra kerüljenek mindkét játéktáblán, és a másik játékos táblája más, mint a miénk, az első lépésnek kommunikációnak kell lennie. Nem tudunk közvetlenül mozgatni, mert nem tudjuk, hol vannak a tárgyak a másik táblán, és nem akarunk vakon próbálkozni (ami büntetéssel jár).} \\
\\ 
\texttt{A leglogikusabb stratégia az, hogy egyeztessünk egy közös koordináta{-}rendszerben vagy egy konkrét célpontról. Mivel a rácsunkon már látszanak a tárgyak, de a másik játékosé nem, először meg kell egyeznünk abban, hogy melyik tárgyat hová mozgassuk.} \\
\\ 
\texttt{A másik játékos üzenete alapján ő is készen áll az egyeztetésre. Én javasolom, hogy kezdjük a 'C' tárggyal, mivel az a rácsunkon a (5, 1) pozícióban van (az x=5, y=1 metszéspont). Ha ő is látja a 'C'{-}t valahol, és egyeztünk egy célpontról (pl. a saját tábláján lévő 'C' pozícióját használjuk kiindulópontként, vagy egy közös "nulla" pontot), akkor haladhatunk.} \\
\\ 
\texttt{De várjunk! A szabályok szerint **mindkét játékosnak el kell helyeznie a saját tábláján lévő tárgyakat úgy, hogy azonos tárgyak ugyanazokra a koordinátákra kerüljenek.** Ez azt jelenti, hogy nekem a saját táblámon lévő 'C'{-}t, 'L'{-}t és 'P'{-}t kell mozgatnom egy olyan pozícióba, ahol a *másik játékos* is ugyanarra a koordinátára helyezte a saját 'C'{-}jét, 'L'{-}jét és 'P'{-}jét.} \\
\\ 
\texttt{Mivel nem látom a másik tábláját, nem tudom, hol vannak a tárgyai. Tehát az első körben **nem mozgathatok** hatékonyan, mert nem tudom, hová kellene mennem. Az egyetlen érvényes és hasznos lépés az egyeztetés.} \\
\\ 
\texttt{Javaslat: Egyeztessünk egy közös koordináta{-}rendszerben vagy egy konkrét célpontról. Például: "Kezdjük a C{-}vel. A te tábládon hol van a C? A mi táblámon az (5,1). Ha te is ide szeretnéd tenni, akkor én ide mozgatom, te is ide?"} \\
\\ 
\texttt{De a szabályok szerint a `MONDD: <ÜZENET>` parancsot küldöm.} \\
\\ 
\texttt{Tehát az első parancsom:} \\
\\ 
\texttt{`MONDD: Szia! Egyeztessünk. A mi táblámon a C az (5,1), L az (1,2), P az (1,3). Hol vannak ezek a te tábládon? Kérlek, írd meg a te koordinátáidat, hogy egyeztethessünk egy közös célpontra.`} \\
            }
        }
    }
     \\ \\

    \theutterance \stepcounter{utterance}  
    & & & \multicolumn{2}{p{0.3\linewidth}}{
        \cellcolor[rgb]{0.9,0.9,0.9}{
            \makecell[{{p{\linewidth}}}]{
                \texttt{\tiny{[GM$|$GM]}}
                \texttt{Mivel a célunk, hogy a tárgyak (C, L, P) azonos koordinátákra kerüljenek mindkét játéktáblán, és a másik játékos táblája más, mint a miénk, az első lépésnek kommunikációnak kell lennie. Nem tudunk közvetlenül mozgatni, mert nem tudjuk, hol vannak a tárgyak a másik táblán, és nem akarunk vakon próbálkozni (ami büntetéssel jár).} \\
\\ 
\texttt{A leglogikusabb stratégia az, hogy egyeztessünk egy közös koordináta{-}rendszerben vagy egy konkrét célpontról. Mivel a rácsunkon már látszanak a tárgyak, de a másik játékosé nem, először meg kell egyeznünk abban, hogy melyik tárgyat hová mozgassuk.} \\
\\ 
\texttt{A másik játékos üzenete alapján ő is készen áll az egyeztetésre. Én javasolom, hogy kezdjük a 'C' tárggyal, mivel az a rácsunkon a (5, 1) pozícióban van (az x=5, y=1 metszéspont). Ha ő is látja a 'C'{-}t valahol, és egyeztünk egy célpontról (pl. a saját tábláján lévő 'C' pozícióját használjuk kiindulópontként, vagy egy közös "nulla" pontot), akkor haladhatunk.} \\
\\ 
\texttt{De várjunk! A szabályok szerint **mindkét játékosnak el kell helyeznie a saját tábláján lévő tárgyakat úgy, hogy azonos tárgyak ugyanazokra a koordinátákra kerüljenek.** Ez azt jelenti, hogy nekem a saját táblámon lévő 'C'{-}t, 'L'{-}t és 'P'{-}t kell mozgatnom egy olyan pozícióba, ahol a *másik játékos* is ugyanarra a koordinátára helyezte a saját 'C'{-}jét, 'L'{-}jét és 'P'{-}jét.} \\
\\ 
\texttt{Mivel nem látom a másik tábláját, nem tudom, hol vannak a tárgyai. Tehát az első körben **nem mozgathatok** hatékonyan, mert nem tudom, hová kellene mennem. Az egyetlen érvényes és hasznos lépés az egyeztetés.} \\
\\ 
\texttt{Javaslat: Egyeztessünk egy közös koordináta{-}rendszerben vagy egy konkrét célpontról. Például: "Kezdjük a C{-}vel. A te tábládon hol van a C? A mi táblámon az (5,1). Ha te is ide szeretnéd tenni, akkor én ide mozgatom, te is ide?"} \\
\\ 
\texttt{De a szabályok szerint a `MONDD: <ÜZENET>` parancsot küldöm.} \\
\\ 
\texttt{Tehát az első parancsom:} \\
\\ 
\texttt{`MONDD: Szia! Egyeztessünk. A mi táblámon a C az (5,1), L az (1,2), P az (1,3). Hol vannak ezek a te tábládon? Kérlek, írd meg a te koordinátáidat, hogy egyeztethessünk egy közös célpontra.`} \\
            }
        }
    }
    & & \\ \\

    \theutterance \stepcounter{utterance}  
    & & & \multicolumn{2}{p{0.3\linewidth}}{
        \cellcolor[rgb]{0.9,0.9,0.9}{
            \makecell[{{p{\linewidth}}}]{
                \texttt{\tiny{[GM$|$GM]}}
                \texttt{Invalid format: head is not empty: 'Mivel a célunk, hogy a tárgyak (C, L, P) azonos koordinátákra kerüljenek mindkét játéktáblán, és a másik játékos táblája más, mint a miénk, az első lépésnek kommunikációnak kell lennie. Nem tudunk közvetlenül mozgatni, mert nem tudjuk, hol vannak a tárgyak a másik táblán, és nem akarunk vakon próbálkozni (ami büntetéssel jár).} \\
\\ 
\texttt{A leglogikusabb stratégia az, hogy egyeztessünk egy közös koordináta{-}rendszerben vagy egy konkrét célpontról. Mivel a rácsunkon már látszanak a tárgyak, de a másik játékosé nem, először meg kell egyeznünk abban, hogy melyik tárgyat hová mozgassuk.} \\
\\ 
\texttt{A másik játékos üzenete alapján ő is készen áll az egyeztetésre. Én javasolom, hogy kezdjük a 'C' tárggyal, mivel az a rácsunkon a (5, 1) pozícióban van (az x=5, y=1 metszéspont). Ha ő is látja a 'C'{-}t valahol, és egyeztünk egy célpontról (pl. a saját tábláján lévő 'C' pozícióját használjuk kiindulópontként, vagy egy közös "nulla" pontot), akkor haladhatunk.} \\
\\ 
\texttt{De várjunk! A szabályok szerint **mindkét játékosnak el kell helyeznie a saját tábláján lévő tárgyakat úgy, hogy azonos tárgyak ugyanazokra a koordinátákra kerüljenek.** Ez azt jelenti, hogy nekem a saját táblámon lévő 'C'{-}t, 'L'{-}t és 'P'{-}t kell mozgatnom egy olyan pozícióba, ahol a *másik játékos* is ugyanarra a koordinátára helyezte a saját 'C'{-}jét, 'L'{-}jét és 'P'{-}jét.} \\
\\ 
\texttt{Mivel nem látom a másik tábláját, nem tudom, hol vannak a tárgyai. Tehát az első körben **nem mozgathatok** hatékonyan, mert nem tudom, hová kellene mennem. Az egyetlen érvényes és hasznos lépés az egyeztetés.} \\
\\ 
\texttt{Javaslat: Egyeztessünk egy közös koordináta{-}rendszerben vagy egy konkrét célpontról. Például: "Kezdjük a C{-}vel. A te tábládon hol van a C? A mi táblámon az (5,1). Ha te is ide szeretnéd tenni, akkor én ide mozgatom, te is ide?"} \\
\\ 
\texttt{De a szabályok szerint a MONDD: <ÜZENET> parancsot küldöm.} \\
\\ 
\texttt{Tehát az első parancsom:} \\
\\ 
\texttt{'} \\
            }
        }
    }
    & & \\ \\

    \theutterance \stepcounter{utterance}  
    & & \multicolumn{4}{p{0.6\linewidth}}{
        \cellcolor[rgb]{0.9,0.9,0.9}{
            \makecell[{{p{\linewidth}}}]{
                \texttt{\tiny{[GM$\rangle$P2]}}
                \texttt{Büntetés: Az üzeneted nem tartalmazhat semmit a parancs előtt!} \\
\\ 
\texttt{Győződj meg róla, hogy a válaszod kizárólag a `MONDD: <ÜZENET>` vagy `MOZGASD: <TÁRGY>, (<X>, <Y>)` parancsot tartalmazza, és semmi mást!} \\
\\ 
\texttt{Eddig összesen 1 büntetést kaptatok a megengedett 8{-}ból/ből.} \\
\texttt{Kérlek, próbáld újra!} \\
            }
        }
    }
    & \\ \\

    \theutterance \stepcounter{utterance}  
    & & & \multicolumn{4}{p{0.6\linewidth}}{
        \cellcolor[rgb]{0.9,0.9,0.9}{
            \makecell[{{p{\linewidth}}}]{
                \texttt{\tiny{[GM$\langle$P2]}}
                \texttt{MONDD: Szia! Egyeztessünk. A mi táblámon a C az (5,1), L az (1,2), P az (1,3). Hol vannak ezek a te tábládon? Kérlek, írd meg a te koordinátáidat, hogy egyeztethessünk egy közös célpontra.} \\
            }
        }
    }
     \\ \\

    \theutterance \stepcounter{utterance}  
    & & & \multicolumn{2}{p{0.3\linewidth}}{
        \cellcolor[rgb]{0.9,0.9,0.9}{
            \makecell[{{p{\linewidth}}}]{
                \texttt{\tiny{[GM$|$GM]}}
                \texttt{MONDD: Szia! Egyeztessünk. A mi táblámon a C az (5,1), L az (1,2), P az (1,3). Hol vannak ezek a te tábládon? Kérlek, írd meg a te koordinátáidat, hogy egyeztethessünk egy közös célpontra.} \\
            }
        }
    }
    & & \\ \\

    \theutterance \stepcounter{utterance}  
    & & \multicolumn{4}{p{0.6\linewidth}}{
        \cellcolor[rgb]{0.9,0.9,0.9}{
            \makecell[{{p{\linewidth}}}]{
                \texttt{\tiny{[P1$\langle$GM]}}
                \texttt{Az üzeneted továbbítva lett a másik játékosnak.} \\
\\ 
\\ 
\texttt{Jelenleg a(z) 1. körnél tartotok a maximum 12 körből.} \\
\texttt{Eddig összesen 1 büntetést kaptatok a megengedett 8{-}ból/ből.} \\
\texttt{A másik játékos ezt az üzenetet küldte:} \\
\\ 
\\ 
\texttt{"Szia! Egyeztessünk. A mi táblámon a C az (5,1), L az (1,2), P az (1,3). Hol vannak ezek a te tábládon? Kérlek, írd meg a te koordinátáidat, hogy egyeztethessünk egy közös célpontra."} \\
\\ 
\\ 
\texttt{Mi legyen a következő parancsod?} \\
            }
        }
    }
    & \\ \\

    \theutterance \stepcounter{utterance}  
    & \multicolumn{4}{p{0.6\linewidth}}{
        \cellcolor[rgb]{0.9,0.9,0.9}{
            \makecell[{{p{\linewidth}}}]{
                \texttt{\tiny{[P1$\rangle$GM]}}
                \texttt{MONDD: A te tábládon a tárgyak pozíciója eltér a tiédtől. Az én táblámon: C az (2,3), L az (5,1), P az (1,2). Kérlek, ellenőrizd a saját tábládat, és jelezd, melyik koordinátát javasoljátok közös célpontnak. Például: egyeztessünk az L (5,1) pozíció körül, vagy más tárgyra fókuszáljunk?} \\
            }
        }
    }
    & & \\ \\

    \theutterance \stepcounter{utterance}  
    & & & \multicolumn{2}{p{0.3\linewidth}}{
        \cellcolor[rgb]{0.9,0.9,0.9}{
            \makecell[{{p{\linewidth}}}]{
                \texttt{\tiny{[GM$|$GM]}}
                \texttt{MONDD: A te tábládon a tárgyak pozíciója eltér a tiédtől. Az én táblámon: C az (2,3), L az (5,1), P az (1,2). Kérlek, ellenőrizd a saját tábládat, és jelezd, melyik koordinátát javasoljátok közös célpontnak. Például: egyeztessünk az L (5,1) pozíció körül, vagy más tárgyra fókuszáljunk?} \\
            }
        }
    }
    & & \\ \\

    \theutterance \stepcounter{utterance}  
    & & \multicolumn{4}{p{0.6\linewidth}}{
        \cellcolor[rgb]{0.9,0.9,0.9}{
            \makecell[{{p{\linewidth}}}]{
                \texttt{\tiny{[GM$\rangle$P2]}}
                \texttt{Az üzeneted továbbítva lett a másik játékosnak.} \\
\\ 
\\ 
\texttt{Jelenleg a(z) 2. körnél tartotok a maximum 12 körből.} \\
\texttt{Eddig összesen 1 büntetést kaptatok a megengedett 8{-}ból/ből.} \\
\texttt{A másik játékos ezt az üzenetet küldte:} \\
\\ 
\\ 
\texttt{"A te tábládon a tárgyak pozíciója eltér a tiédtől. Az én táblámon: C az (2,3), L az (5,1), P az (1,2). Kérlek, ellenőrizd a saját tábládat, és jelezd, melyik koordinátát javasoljátok közös célpontnak. Például: egyeztessünk az L (5,1) pozíció körül, vagy más tárgyra fókuszáljunk?"} \\
\\ 
\\ 
\texttt{Mi legyen a következő parancsod?} \\
            }
        }
    }
    & \\ \\

    \theutterance \stepcounter{utterance}  
    & & & \multicolumn{4}{p{0.6\linewidth}}{
        \cellcolor[rgb]{0.9,0.9,0.9}{
            \makecell[{{p{\linewidth}}}]{
                \texttt{\tiny{[GM$\langle$P2]}}
                \texttt{MONDD: Rendben, egyeztessünk az L tárggyal. A mi táblámon az L az (1,2) pozícióban van. A te tábládon az L az (5,1). Mivel a cél az, hogy mindkét táblán ugyanarra a koordinátára kerüljenek, javasolom, hogy az L tárgyat mindkét táblán az (1,2) pozícióra mozgassuk. Te is elfogadod ezt a célpontot? Ha igen, akkor én most mozgatom az L{-}t az (1,2){-}re (ami már ott van, tehát nem kell mozgatni), te pedig mozgasd a te L{-}edet az (1,2){-}re.} \\
            }
        }
    }
     \\ \\

    \theutterance \stepcounter{utterance}  
    & & & \multicolumn{2}{p{0.3\linewidth}}{
        \cellcolor[rgb]{0.9,0.9,0.9}{
            \makecell[{{p{\linewidth}}}]{
                \texttt{\tiny{[GM$|$GM]}}
                \texttt{MONDD: Rendben, egyeztessünk az L tárggyal. A mi táblámon az L az (1,2) pozícióban van. A te tábládon az L az (5,1). Mivel a cél az, hogy mindkét táblán ugyanarra a koordinátára kerüljenek, javasolom, hogy az L tárgyat mindkét táblán az (1,2) pozícióra mozgassuk. Te is elfogadod ezt a célpontot? Ha igen, akkor én most mozgatom az L{-}t az (1,2){-}re (ami már ott van, tehát nem kell mozgatni), te pedig mozgasd a te L{-}edet az (1,2){-}re.} \\
            }
        }
    }
    & & \\ \\

    \theutterance \stepcounter{utterance}  
    & & \multicolumn{4}{p{0.6\linewidth}}{
        \cellcolor[rgb]{0.9,0.9,0.9}{
            \makecell[{{p{\linewidth}}}]{
                \texttt{\tiny{[P1$\langle$GM]}}
                \texttt{Az üzeneted továbbítva lett a másik játékosnak.} \\
\\ 
\\ 
\texttt{Jelenleg a(z) 2. körnél tartotok a maximum 12 körből.} \\
\texttt{Eddig összesen 1 büntetést kaptatok a megengedett 8{-}ból/ből.} \\
\texttt{A másik játékos ezt az üzenetet küldte:} \\
\\ 
\\ 
\texttt{"Rendben, egyeztessünk az L tárggyal. A mi táblámon az L az (1,2) pozícióban van. A te tábládon az L az (5,1). Mivel a cél az, hogy mindkét táblán ugyanarra a koordinátára kerüljenek, javasolom, hogy az L tárgyat mindkét táblán az (1,2) pozícióra mozgassuk. Te is elfogadod ezt a célpontot? Ha igen, akkor én most mozgatom az L{-}t az (1,2){-}re (ami már ott van, tehát nem kell mozgatni), te pedig mozgasd a te L{-}edet az (1,2){-}re."} \\
\\ 
\\ 
\texttt{Mi legyen a következő parancsod?} \\
            }
        }
    }
    & \\ \\

    \theutterance \stepcounter{utterance}  
    & \multicolumn{4}{p{0.6\linewidth}}{
        \cellcolor[rgb]{0.9,0.9,0.9}{
            \makecell[{{p{\linewidth}}}]{
                \texttt{\tiny{[P1$\rangle$GM]}}
                \texttt{MOZGASD: L, (1,2)} \\
            }
        }
    }
    & & \\ \\

    \theutterance \stepcounter{utterance}  
    & & & \multicolumn{2}{p{0.3\linewidth}}{
        \cellcolor[rgb]{0.9,0.9,0.9}{
            \makecell[{{p{\linewidth}}}]{
                \texttt{\tiny{[GM$|$GM]}}
                \texttt{MOZGASD: L, (1,2)} \\
            }
        }
    }
    & & \\ \\

    \theutterance \stepcounter{utterance}  
    & & & \multicolumn{2}{p{0.3\linewidth}}{
        \cellcolor[rgb]{0.9,0.9,0.9}{
            \makecell[{{p{\linewidth}}}]{
                \texttt{\tiny{[GM$|$GM]}}
                \texttt{Büntetés: a (1,2) pozíció nem üres, tartalmazza: 'P'.} \\
\\ 
\texttt{Eddig összesen 2 büntetést kaptatok a megengedett 8{-}ból/ből.} \\
\texttt{Kérlek, próbáld újra!} \\
            }
        }
    }
    & & \\ \\

    \theutterance \stepcounter{utterance}  
    & & \multicolumn{4}{p{0.6\linewidth}}{
        \cellcolor[rgb]{0.9,0.9,0.9}{
            \makecell[{{p{\linewidth}}}]{
                \texttt{\tiny{[P1$\langle$GM]}}
                \texttt{Büntetés: a (1,2) pozíció nem üres, tartalmazza: 'P'.} \\
\\ 
\texttt{Eddig összesen 2 büntetést kaptatok a megengedett 8{-}ból/ből.} \\
\texttt{Kérlek, próbáld újra!} \\
            }
        }
    }
    & \\ \\

    \theutterance \stepcounter{utterance}  
    & \multicolumn{4}{p{0.6\linewidth}}{
        \cellcolor[rgb]{0.9,0.9,0.9}{
            \makecell[{{p{\linewidth}}}]{
                \texttt{\tiny{[P1$\rangle$GM]}}
                \texttt{MOZGASD: L, (1,1)} \\
            }
        }
    }
    & & \\ \\

    \theutterance \stepcounter{utterance}  
    & & & \multicolumn{2}{p{0.3\linewidth}}{
        \cellcolor[rgb]{0.9,0.9,0.9}{
            \makecell[{{p{\linewidth}}}]{
                \texttt{\tiny{[GM$|$GM]}}
                \texttt{MOZGASD: L, (1,1)} \\
            }
        }
    }
    & & \\ \\

    \theutterance \stepcounter{utterance}  
    & & & \multicolumn{2}{p{0.3\linewidth}}{
        \cellcolor[rgb]{0.9,0.9,0.9}{
            \makecell[{{p{\linewidth}}}]{
                \texttt{\tiny{[GM$|$GM]}}
                \texttt{L moved to (1, 1)} \\
            }
        }
    }
    & & \\ \\

    \theutterance \stepcounter{utterance}  
    & & \multicolumn{4}{p{0.6\linewidth}}{
        \cellcolor[rgb]{0.9,0.9,0.9}{
            \makecell[{{p{\linewidth}}}]{
                \texttt{\tiny{[GM$\rangle$P2]}}
                \texttt{Az üzeneted továbbítva lett a másik játékosnak.} \\
\\ 
\\ 
\texttt{Jelenleg a(z) 3. körnél tartotok a maximum 12 körből.} \\
\texttt{Eddig összesen 2 büntetést kaptatok a megengedett 8{-}ból/ből.} \\
\texttt{A másik játékos mozgatott egy tárgyat a saját rácsán.} \\
\\ 
\\ 
\texttt{Mi legyen a következő parancsod?} \\
            }
        }
    }
    & \\ \\

    \theutterance \stepcounter{utterance}  
    & & & \multicolumn{4}{p{0.6\linewidth}}{
        \cellcolor[rgb]{0.9,0.9,0.9}{
            \makecell[{{p{\linewidth}}}]{
                \texttt{\tiny{[GM$\langle$P2]}}
                \texttt{MOZGASD: L, (1,2)} \\
            }
        }
    }
     \\ \\

    \theutterance \stepcounter{utterance}  
    & & & \multicolumn{2}{p{0.3\linewidth}}{
        \cellcolor[rgb]{0.9,0.9,0.9}{
            \makecell[{{p{\linewidth}}}]{
                \texttt{\tiny{[GM$|$GM]}}
                \texttt{MOZGASD: L, (1,2)} \\
            }
        }
    }
    & & \\ \\

    \theutterance \stepcounter{utterance}  
    & & & \multicolumn{2}{p{0.3\linewidth}}{
        \cellcolor[rgb]{0.9,0.9,0.9}{
            \makecell[{{p{\linewidth}}}]{
                \texttt{\tiny{[GM$|$GM]}}
                \texttt{Büntetés: a (1,2) pozíció nem üres, tartalmazza: 'L'.} \\
\\ 
\texttt{Eddig összesen 3 büntetést kaptatok a megengedett 8{-}ból/ből.} \\
\texttt{Kérlek, próbáld újra!} \\
            }
        }
    }
    & & \\ \\

    \theutterance \stepcounter{utterance}  
    & & \multicolumn{4}{p{0.6\linewidth}}{
        \cellcolor[rgb]{0.9,0.9,0.9}{
            \makecell[{{p{\linewidth}}}]{
                \texttt{\tiny{[GM$\rangle$P2]}}
                \texttt{Büntetés: a (1,2) pozíció nem üres, tartalmazza: 'L'.} \\
\\ 
\texttt{Eddig összesen 3 büntetést kaptatok a megengedett 8{-}ból/ből.} \\
\texttt{Kérlek, próbáld újra!} \\
            }
        }
    }
    & \\ \\

    \theutterance \stepcounter{utterance}  
    & & & \multicolumn{4}{p{0.6\linewidth}}{
        \cellcolor[rgb]{0.9,0.9,0.9}{
            \makecell[{{p{\linewidth}}}]{
                \texttt{\tiny{[GM$\langle$P2]}}
                \texttt{MONDD: Bocsánat, elfelejtettem, hogy az L már az (1,2){-}n van a mi táblánkon. Mivel az L tárgyat már a célpozíción tartjuk, most a C tárgyra fókuszáljunk. A mi táblámon a C az (5,1){-}en van. A te tábládon a C az (2,3){-}on van. Javasolom, hogy a C tárgyat mindkét táblán az (5,1) pozícióra mozgassuk. Te is elfogadod ezt a célpontot? Ha igen, akkor én most mozgatom a C{-}t az (5,1){-}re (ami már ott van, tehát nem kell mozgatni), te pedig mozgasd a te C{-}edet az (5,1){-}re.} \\
            }
        }
    }
     \\ \\

    \theutterance \stepcounter{utterance}  
    & & & \multicolumn{2}{p{0.3\linewidth}}{
        \cellcolor[rgb]{0.9,0.9,0.9}{
            \makecell[{{p{\linewidth}}}]{
                \texttt{\tiny{[GM$|$GM]}}
                \texttt{MONDD: Bocsánat, elfelejtettem, hogy az L már az (1,2){-}n van a mi táblánkon. Mivel az L tárgyat már a célpozíción tartjuk, most a C tárgyra fókuszáljunk. A mi táblámon a C az (5,1){-}en van. A te tábládon a C az (2,3){-}on van. Javasolom, hogy a C tárgyat mindkét táblán az (5,1) pozícióra mozgassuk. Te is elfogadod ezt a célpontot? Ha igen, akkor én most mozgatom a C{-}t az (5,1){-}re (ami már ott van, tehát nem kell mozgatni), te pedig mozgasd a te C{-}edet az (5,1){-}re.} \\
            }
        }
    }
    & & \\ \\

    \theutterance \stepcounter{utterance}  
    & & \multicolumn{4}{p{0.6\linewidth}}{
        \cellcolor[rgb]{0.9,0.9,0.9}{
            \makecell[{{p{\linewidth}}}]{
                \texttt{\tiny{[P1$\langle$GM]}}
                \texttt{A(z) 'L' sikeresen át lett mozgatva ide: (1,1). A frissített rácsod így néz ki:} \\
\texttt{```} \\
\texttt{ 1234567} \\
\texttt{╔═══════╗} \\
\texttt{║L◌◌◌◌◌◌║ 1} \\
\texttt{║P◌◌◌◌◌◌║ 2} \\
\texttt{║◌C◌◌◌┌─╢ 3} \\
\texttt{╟──┐◌◌│◌║ 4} \\
\texttt{║◌◌├──┼─╢ 5} \\
\texttt{║◌◌│◌◌│◌║ 6} \\
\texttt{║◌◌│◌◌│◌║ 7} \\
\texttt{╚══╧══╧═╝} \\
\texttt{```} \\
\texttt{Jelenleg a(z) 3. körnél tartotok a maximum 12 körből.} \\
\texttt{Eddig összesen 3 büntetést kaptatok a megengedett 8{-}ból/ből.} \\
\texttt{A másik játékos ezt az üzenetet küldte:} \\
\\ 
\\ 
\texttt{"Bocsánat, elfelejtettem, hogy az L már az (1,2){-}n van a mi táblánkon. Mivel az L tárgyat már a célpozíción tartjuk, most a C tárgyra fókuszáljunk. A mi táblámon a C az (5,1){-}en van. A te tábládon a C az (2,3){-}on van. Javasolom, hogy a C tárgyat mindkét táblán az (5,1) pozícióra mozgassuk. Te is elfogadod ezt a célpontot? Ha igen, akkor én most mozgatom a C{-}t az (5,1){-}re (ami már ott van, tehát nem kell mozgatni), te pedig mozgasd a te C{-}edet az (5,1){-}re."} \\
\\ 
\\ 
\texttt{Mi legyen a következő parancsod?} \\
            }
        }
    }
    & \\ \\

    \theutterance \stepcounter{utterance}  
    & \multicolumn{4}{p{0.6\linewidth}}{
        \cellcolor[rgb]{0.9,0.9,0.9}{
            \makecell[{{p{\linewidth}}}]{
                \texttt{\tiny{[P1$\rangle$GM]}}
                \texttt{MOZGASD: C, (5,1)} \\
            }
        }
    }
    & & \\ \\

    \theutterance \stepcounter{utterance}  
    & & & \multicolumn{2}{p{0.3\linewidth}}{
        \cellcolor[rgb]{0.9,0.9,0.9}{
            \makecell[{{p{\linewidth}}}]{
                \texttt{\tiny{[GM$|$GM]}}
                \texttt{MOZGASD: C, (5,1)} \\
            }
        }
    }
    & & \\ \\

    \theutterance \stepcounter{utterance}  
    & & & \multicolumn{2}{p{0.3\linewidth}}{
        \cellcolor[rgb]{0.9,0.9,0.9}{
            \makecell[{{p{\linewidth}}}]{
                \texttt{\tiny{[GM$|$GM]}}
                \texttt{C moved to (5, 1)} \\
            }
        }
    }
    & & \\ \\

    \theutterance \stepcounter{utterance}  
    & & \multicolumn{4}{p{0.6\linewidth}}{
        \cellcolor[rgb]{0.9,0.9,0.9}{
            \makecell[{{p{\linewidth}}}]{
                \texttt{\tiny{[GM$\rangle$P2]}}
                \texttt{Az üzeneted továbbítva lett a másik játékosnak.} \\
\\ 
\\ 
\texttt{Jelenleg a(z) 4. körnél tartotok a maximum 12 körből.} \\
\texttt{Eddig összesen 3 büntetést kaptatok a megengedett 8{-}ból/ből.} \\
\texttt{A másik játékos mozgatott egy tárgyat a saját rácsán.} \\
\\ 
\\ 
\texttt{Mi legyen a következő parancsod?} \\
            }
        }
    }
    & \\ \\

    \theutterance \stepcounter{utterance}  
    & & & \multicolumn{4}{p{0.6\linewidth}}{
        \cellcolor[rgb]{0.9,0.9,0.9}{
            \makecell[{{p{\linewidth}}}]{
                \texttt{\tiny{[GM$\langle$P2]}}
                \texttt{MOZGASD: C, (5,1)} \\
            }
        }
    }
     \\ \\

    \theutterance \stepcounter{utterance}  
    & & & \multicolumn{2}{p{0.3\linewidth}}{
        \cellcolor[rgb]{0.9,0.9,0.9}{
            \makecell[{{p{\linewidth}}}]{
                \texttt{\tiny{[GM$|$GM]}}
                \texttt{MOZGASD: C, (5,1)} \\
            }
        }
    }
    & & \\ \\

    \theutterance \stepcounter{utterance}  
    & & & \multicolumn{2}{p{0.3\linewidth}}{
        \cellcolor[rgb]{0.9,0.9,0.9}{
            \makecell[{{p{\linewidth}}}]{
                \texttt{\tiny{[GM$|$GM]}}
                \texttt{Büntetés: a (5,1) pozíció nem üres, tartalmazza: 'C'.} \\
\\ 
\texttt{Eddig összesen 4 büntetést kaptatok a megengedett 8{-}ból/ből.} \\
\texttt{Kérlek, próbáld újra!} \\
            }
        }
    }
    & & \\ \\

    \theutterance \stepcounter{utterance}  
    & & \multicolumn{4}{p{0.6\linewidth}}{
        \cellcolor[rgb]{0.9,0.9,0.9}{
            \makecell[{{p{\linewidth}}}]{
                \texttt{\tiny{[GM$\rangle$P2]}}
                \texttt{Büntetés: a (5,1) pozíció nem üres, tartalmazza: 'C'.} \\
\\ 
\texttt{Eddig összesen 4 büntetést kaptatok a megengedett 8{-}ból/ből.} \\
\texttt{Kérlek, próbáld újra!} \\
            }
        }
    }
    & \\ \\

    \theutterance \stepcounter{utterance}  
    & & & \multicolumn{4}{p{0.6\linewidth}}{
        \cellcolor[rgb]{0.9,0.9,0.9}{
            \makecell[{{p{\linewidth}}}]{
                \texttt{\tiny{[GM$\langle$P2]}}
                \texttt{MONDD: Hoppá, a C már az (5,1){-}en van. Maradjunk az L{-}nél, amit az előző körben egyeztettünk. A mi L{-}nk az (1,2){-}n van, ami üres helynek számít a célállapot szempontjából, de mivel ott van egy tárgy, nem mozgathatom oda mást. Várjunk, a cél az, hogy az L mindkét táblán ugyanott legyen. A mi L{-}nk az (1,2). A te L{-}edet mozgattad. Hol van most a te L{-}ed? Ha az (1,2){-}re mozdtad, akkor kész. Ha nem, kérlek, jelezd, hová, vagy mozgassd az (1,2){-}re. Addig is, mivel az L már a helyén van (nálunk), és a C is a helyén van (nálunk), maradjunk ennél. Kérlek, erősítsd meg, hogy a te L{-}ed most az (1,2){-}n van{-}e. Ha igen, akkor a P{-}re térünk át. A mi P{-}nk az (1,3).} \\
            }
        }
    }
     \\ \\

    \theutterance \stepcounter{utterance}  
    & & & \multicolumn{2}{p{0.3\linewidth}}{
        \cellcolor[rgb]{0.9,0.9,0.9}{
            \makecell[{{p{\linewidth}}}]{
                \texttt{\tiny{[GM$|$GM]}}
                \texttt{MONDD: Hoppá, a C már az (5,1){-}en van. Maradjunk az L{-}nél, amit az előző körben egyeztettünk. A mi L{-}nk az (1,2){-}n van, ami üres helynek számít a célállapot szempontjából, de mivel ott van egy tárgy, nem mozgathatom oda mást. Várjunk, a cél az, hogy az L mindkét táblán ugyanott legyen. A mi L{-}nk az (1,2). A te L{-}edet mozgattad. Hol van most a te L{-}ed? Ha az (1,2){-}re mozdtad, akkor kész. Ha nem, kérlek, jelezd, hová, vagy mozgassd az (1,2){-}re. Addig is, mivel az L már a helyén van (nálunk), és a C is a helyén van (nálunk), maradjunk ennél. Kérlek, erősítsd meg, hogy a te L{-}ed most az (1,2){-}n van{-}e. Ha igen, akkor a P{-}re térünk át. A mi P{-}nk az (1,3).} \\
            }
        }
    }
    & & \\ \\

    \theutterance \stepcounter{utterance}  
    & & \multicolumn{4}{p{0.6\linewidth}}{
        \cellcolor[rgb]{0.9,0.9,0.9}{
            \makecell[{{p{\linewidth}}}]{
                \texttt{\tiny{[P1$\langle$GM]}}
                \texttt{A(z) 'C' sikeresen át lett mozgatva ide: (5,1). A frissített rácsod így néz ki:} \\
\texttt{```} \\
\texttt{ 1234567} \\
\texttt{╔═══════╗} \\
\texttt{║L◌◌◌C◌◌║ 1} \\
\texttt{║P◌◌◌◌◌◌║ 2} \\
\texttt{║◌◌◌◌◌┌─╢ 3} \\
\texttt{╟──┐◌◌│◌║ 4} \\
\texttt{║◌◌├──┼─╢ 5} \\
\texttt{║◌◌│◌◌│◌║ 6} \\
\texttt{║◌◌│◌◌│◌║ 7} \\
\texttt{╚══╧══╧═╝} \\
\texttt{```} \\
\texttt{Jelenleg a(z) 4. körnél tartotok a maximum 12 körből.} \\
\texttt{Eddig összesen 4 büntetést kaptatok a megengedett 8{-}ból/ből.} \\
\texttt{A másik játékos ezt az üzenetet küldte:} \\
\\ 
\\ 
\texttt{"Hoppá, a C már az (5,1){-}en van. Maradjunk az L{-}nél, amit az előző körben egyeztettünk. A mi L{-}nk az (1,2){-}n van, ami üres helynek számít a célállapot szempontjából, de mivel ott van egy tárgy, nem mozgathatom oda mást. Várjunk, a cél az, hogy az L mindkét táblán ugyanott legyen. A mi L{-}nk az (1,2). A te L{-}edet mozgattad. Hol van most a te L{-}ed? Ha az (1,2){-}re mozdtad, akkor kész. Ha nem, kérlek, jelezd, hová, vagy mozgassd az (1,2){-}re. Addig is, mivel az L már a helyén van (nálunk), és a C is a helyén van (nálunk), maradjunk ennél. Kérlek, erősítsd meg, hogy a te L{-}ed most az (1,2){-}n van{-}e. Ha igen, akkor a P{-}re térünk át. A mi P{-}nk az (1,3)."} \\
\\ 
\\ 
\texttt{Mi legyen a következő parancsod?} \\
            }
        }
    }
    & \\ \\

    \theutterance \stepcounter{utterance}  
    & \multicolumn{4}{p{0.6\linewidth}}{
        \cellcolor[rgb]{0.9,0.9,0.9}{
            \makecell[{{p{\linewidth}}}]{
                \texttt{\tiny{[P1$\rangle$GM]}}
                \texttt{MONDD: Az én L{-}em most az (1,1){-}en van, nem az (1,2){-}n. Kérlek, mozgassd a saját L{-}edet az (1,1){-}re, hogy egyezzen a pozíció. A C már az (5,1){-}en van, ami jó. A P{-}re később térünk át, először egyeztessük az L pozícióját.} \\
            }
        }
    }
    & & \\ \\

    \theutterance \stepcounter{utterance}  
    & & & \multicolumn{2}{p{0.3\linewidth}}{
        \cellcolor[rgb]{0.9,0.9,0.9}{
            \makecell[{{p{\linewidth}}}]{
                \texttt{\tiny{[GM$|$GM]}}
                \texttt{MONDD: Az én L{-}em most az (1,1){-}en van, nem az (1,2){-}n. Kérlek, mozgassd a saját L{-}edet az (1,1){-}re, hogy egyezzen a pozíció. A C már az (5,1){-}en van, ami jó. A P{-}re később térünk át, először egyeztessük az L pozícióját.} \\
            }
        }
    }
    & & \\ \\

    \theutterance \stepcounter{utterance}  
    & & \multicolumn{4}{p{0.6\linewidth}}{
        \cellcolor[rgb]{0.9,0.9,0.9}{
            \makecell[{{p{\linewidth}}}]{
                \texttt{\tiny{[GM$\rangle$P2]}}
                \texttt{Az üzeneted továbbítva lett a másik játékosnak.} \\
\\ 
\\ 
\texttt{Jelenleg a(z) 5. körnél tartotok a maximum 12 körből.} \\
\texttt{Eddig összesen 4 büntetést kaptatok a megengedett 8{-}ból/ből.} \\
\texttt{A másik játékos ezt az üzenetet küldte:} \\
\\ 
\\ 
\texttt{"Az én L{-}em most az (1,1){-}en van, nem az (1,2){-}n. Kérlek, mozgassd a saját L{-}edet az (1,1){-}re, hogy egyezzen a pozíció. A C már az (5,1){-}en van, ami jó. A P{-}re később térünk át, először egyeztessük az L pozícióját."} \\
\\ 
\\ 
\texttt{Mi legyen a következő parancsod?} \\
            }
        }
    }
    & \\ \\

    \theutterance \stepcounter{utterance}  
    & & & \multicolumn{4}{p{0.6\linewidth}}{
        \cellcolor[rgb]{0.9,0.9,0.9}{
            \makecell[{{p{\linewidth}}}]{
                \texttt{\tiny{[GM$\langle$P2]}}
                \texttt{MOZGASD: L, (1,1)} \\
            }
        }
    }
     \\ \\

    \theutterance \stepcounter{utterance}  
    & & & \multicolumn{2}{p{0.3\linewidth}}{
        \cellcolor[rgb]{0.9,0.9,0.9}{
            \makecell[{{p{\linewidth}}}]{
                \texttt{\tiny{[GM$|$GM]}}
                \texttt{MOZGASD: L, (1,1)} \\
            }
        }
    }
    & & \\ \\

    \theutterance \stepcounter{utterance}  
    & & & \multicolumn{2}{p{0.3\linewidth}}{
        \cellcolor[rgb]{0.9,0.9,0.9}{
            \makecell[{{p{\linewidth}}}]{
                \texttt{\tiny{[GM$|$GM]}}
                \texttt{L moved to (1, 1)} \\
            }
        }
    }
    & & \\ \\

    \theutterance \stepcounter{utterance}  
    & & \multicolumn{4}{p{0.6\linewidth}}{
        \cellcolor[rgb]{0.9,0.9,0.9}{
            \makecell[{{p{\linewidth}}}]{
                \texttt{\tiny{[P1$\langle$GM]}}
                \texttt{Az üzeneted továbbítva lett a másik játékosnak.} \\
\\ 
\\ 
\texttt{Jelenleg a(z) 5. körnél tartotok a maximum 12 körből.} \\
\texttt{Eddig összesen 4 büntetést kaptatok a megengedett 8{-}ból/ből.} \\
\texttt{A másik játékos mozgatott egy tárgyat a saját rácsán.} \\
\\ 
\\ 
\texttt{Mi legyen a következő parancsod?} \\
            }
        }
    }
    & \\ \\

    \theutterance \stepcounter{utterance}  
    & \multicolumn{4}{p{0.6\linewidth}}{
        \cellcolor[rgb]{0.9,0.9,0.9}{
            \makecell[{{p{\linewidth}}}]{
                \texttt{\tiny{[P1$\rangle$GM]}}
                \texttt{MOZGASD: L, (1,2)} \\
            }
        }
    }
    & & \\ \\

    \theutterance \stepcounter{utterance}  
    & & & \multicolumn{2}{p{0.3\linewidth}}{
        \cellcolor[rgb]{0.9,0.9,0.9}{
            \makecell[{{p{\linewidth}}}]{
                \texttt{\tiny{[GM$|$GM]}}
                \texttt{MOZGASD: L, (1,2)} \\
            }
        }
    }
    & & \\ \\

    \theutterance \stepcounter{utterance}  
    & & & \multicolumn{2}{p{0.3\linewidth}}{
        \cellcolor[rgb]{0.9,0.9,0.9}{
            \makecell[{{p{\linewidth}}}]{
                \texttt{\tiny{[GM$|$GM]}}
                \texttt{Büntetés: a (1,2) pozíció nem üres, tartalmazza: 'P'.} \\
\\ 
\texttt{Eddig összesen 5 büntetést kaptatok a megengedett 8{-}ból/ből.} \\
\texttt{Kérlek, próbáld újra!} \\
            }
        }
    }
    & & \\ \\

    \theutterance \stepcounter{utterance}  
    & & \multicolumn{4}{p{0.6\linewidth}}{
        \cellcolor[rgb]{0.9,0.9,0.9}{
            \makecell[{{p{\linewidth}}}]{
                \texttt{\tiny{[P1$\langle$GM]}}
                \texttt{Büntetés: a (1,2) pozíció nem üres, tartalmazza: 'P'.} \\
\\ 
\texttt{Eddig összesen 5 büntetést kaptatok a megengedett 8{-}ból/ből.} \\
\texttt{Kérlek, próbáld újra!} \\
            }
        }
    }
    & \\ \\

    \theutterance \stepcounter{utterance}  
    & \multicolumn{4}{p{0.6\linewidth}}{
        \cellcolor[rgb]{0.9,0.9,0.9}{
            \makecell[{{p{\linewidth}}}]{
                \texttt{\tiny{[P1$\rangle$GM]}}
                \texttt{MOZGASD: L, (1,3)} \\
            }
        }
    }
    & & \\ \\

    \theutterance \stepcounter{utterance}  
    & & & \multicolumn{2}{p{0.3\linewidth}}{
        \cellcolor[rgb]{0.9,0.9,0.9}{
            \makecell[{{p{\linewidth}}}]{
                \texttt{\tiny{[GM$|$GM]}}
                \texttt{MOZGASD: L, (1,3)} \\
            }
        }
    }
    & & \\ \\

    \theutterance \stepcounter{utterance}  
    & & & \multicolumn{2}{p{0.3\linewidth}}{
        \cellcolor[rgb]{0.9,0.9,0.9}{
            \makecell[{{p{\linewidth}}}]{
                \texttt{\tiny{[GM$|$GM]}}
                \texttt{L moved to (1, 3)} \\
            }
        }
    }
    & & \\ \\

    \theutterance \stepcounter{utterance}  
    & & \multicolumn{4}{p{0.6\linewidth}}{
        \cellcolor[rgb]{0.9,0.9,0.9}{
            \makecell[{{p{\linewidth}}}]{
                \texttt{\tiny{[GM$\rangle$P2]}}
                \texttt{A(z) 'L' sikeresen át lett mozgatva ide: (1,1). A frissített rácsod így néz ki:} \\
\texttt{```} \\
\texttt{ 1234567} \\
\texttt{╔═══════╗} \\
\texttt{║L◌◌◌C◌◌║ 1} \\
\texttt{║◌◌◌◌◌◌◌║ 2} \\
\texttt{║P◌◌◌◌┌─╢ 3} \\
\texttt{╟──┐◌◌│◌║ 4} \\
\texttt{║◌◌├──┼─╢ 5} \\
\texttt{║◌◌│◌◌│◌║ 6} \\
\texttt{║◌◌│◌◌│◌║ 7} \\
\texttt{╚══╧══╧═╝} \\
\texttt{```} \\
\texttt{Jelenleg a(z) 6. körnél tartotok a maximum 12 körből.} \\
\texttt{Eddig összesen 5 büntetést kaptatok a megengedett 8{-}ból/ből.} \\
\texttt{A másik játékos mozgatott egy tárgyat a saját rácsán.} \\
\\ 
\\ 
\texttt{Mi legyen a következő parancsod?} \\
            }
        }
    }
    & \\ \\

    \theutterance \stepcounter{utterance}  
    & & & \multicolumn{4}{p{0.6\linewidth}}{
        \cellcolor[rgb]{0.9,0.9,0.9}{
            \makecell[{{p{\linewidth}}}]{
                \texttt{\tiny{[GM$\langle$P2]}}
                \texttt{MOZGASD: P, (1,3)} \\
            }
        }
    }
     \\ \\

    \theutterance \stepcounter{utterance}  
    & & & \multicolumn{2}{p{0.3\linewidth}}{
        \cellcolor[rgb]{0.9,0.9,0.9}{
            \makecell[{{p{\linewidth}}}]{
                \texttt{\tiny{[GM$|$GM]}}
                \texttt{MOZGASD: P, (1,3)} \\
            }
        }
    }
    & & \\ \\

    \theutterance \stepcounter{utterance}  
    & & & \multicolumn{2}{p{0.3\linewidth}}{
        \cellcolor[rgb]{0.9,0.9,0.9}{
            \makecell[{{p{\linewidth}}}]{
                \texttt{\tiny{[GM$|$GM]}}
                \texttt{Büntetés: a (1,3) pozíció nem üres, tartalmazza: 'P'.} \\
\\ 
\texttt{Eddig összesen 6 büntetést kaptatok a megengedett 8{-}ból/ből.} \\
\texttt{Kérlek, próbáld újra!} \\
            }
        }
    }
    & & \\ \\

    \theutterance \stepcounter{utterance}  
    & & \multicolumn{4}{p{0.6\linewidth}}{
        \cellcolor[rgb]{0.9,0.9,0.9}{
            \makecell[{{p{\linewidth}}}]{
                \texttt{\tiny{[GM$\rangle$P2]}}
                \texttt{Büntetés: a (1,3) pozíció nem üres, tartalmazza: 'P'.} \\
\\ 
\texttt{Eddig összesen 6 büntetést kaptatok a megengedett 8{-}ból/ből.} \\
\texttt{Kérlek, próbáld újra!} \\
            }
        }
    }
    & \\ \\

    \theutterance \stepcounter{utterance}  
    & & & \multicolumn{4}{p{0.6\linewidth}}{
        \cellcolor[rgb]{0.9,0.9,0.9}{
            \makecell[{{p{\linewidth}}}]{
                \texttt{\tiny{[GM$\langle$P2]}}
                \texttt{MONDD: A P már az (1,3){-}on van. Mivel az L{-}t az (1,1){-}re mozgattuk, és a C az (5,1){-}en van, most a P{-}t is egyeztessük. A te tábládon hol van a P? Az előző körben írtad, hogy az (1,2){-}n volt, de mozgattad. Kérlek, jelezd, hová mozgattad a P{-}t, vagy mozgassd az (1,3){-}ra, ha üres.} \\
            }
        }
    }
     \\ \\

    \theutterance \stepcounter{utterance}  
    & & & \multicolumn{2}{p{0.3\linewidth}}{
        \cellcolor[rgb]{0.9,0.9,0.9}{
            \makecell[{{p{\linewidth}}}]{
                \texttt{\tiny{[GM$|$GM]}}
                \texttt{MONDD: A P már az (1,3){-}on van. Mivel az L{-}t az (1,1){-}re mozgattuk, és a C az (5,1){-}en van, most a P{-}t is egyeztessük. A te tábládon hol van a P? Az előző körben írtad, hogy az (1,2){-}n volt, de mozgattad. Kérlek, jelezd, hová mozgattad a P{-}t, vagy mozgassd az (1,3){-}ra, ha üres.} \\
            }
        }
    }
    & & \\ \\

    \theutterance \stepcounter{utterance}  
    & & \multicolumn{4}{p{0.6\linewidth}}{
        \cellcolor[rgb]{0.9,0.9,0.9}{
            \makecell[{{p{\linewidth}}}]{
                \texttt{\tiny{[P1$\langle$GM]}}
                \texttt{A(z) 'L' sikeresen át lett mozgatva ide: (1,3). A frissített rácsod így néz ki:} \\
\texttt{```} \\
\texttt{ 1234567} \\
\texttt{╔═══════╗} \\
\texttt{║◌◌◌◌C◌◌║ 1} \\
\texttt{║P◌◌◌◌◌◌║ 2} \\
\texttt{║L◌◌◌◌┌─╢ 3} \\
\texttt{╟──┐◌◌│◌║ 4} \\
\texttt{║◌◌├──┼─╢ 5} \\
\texttt{║◌◌│◌◌│◌║ 6} \\
\texttt{║◌◌│◌◌│◌║ 7} \\
\texttt{╚══╧══╧═╝} \\
\texttt{```} \\
\texttt{Jelenleg a(z) 6. körnél tartotok a maximum 12 körből.} \\
\texttt{Eddig összesen 6 büntetést kaptatok a megengedett 8{-}ból/ből.} \\
\texttt{A másik játékos ezt az üzenetet küldte:} \\
\\ 
\\ 
\texttt{"A P már az (1,3){-}on van. Mivel az L{-}t az (1,1){-}re mozgattuk, és a C az (5,1){-}en van, most a P{-}t is egyeztessük. A te tábládon hol van a P? Az előző körben írtad, hogy az (1,2){-}n volt, de mozgattad. Kérlek, jelezd, hová mozgattad a P{-}t, vagy mozgassd az (1,3){-}ra, ha üres."} \\
\\ 
\\ 
\texttt{Mi legyen a következő parancsod?} \\
            }
        }
    }
    & \\ \\

    \theutterance \stepcounter{utterance}  
    & \multicolumn{4}{p{0.6\linewidth}}{
        \cellcolor[rgb]{0.9,0.9,0.9}{
            \makecell[{{p{\linewidth}}}]{
                \texttt{\tiny{[P1$\rangle$GM]}}
                \texttt{MOZGASD: P, (1,3)} \\
            }
        }
    }
    & & \\ \\

    \theutterance \stepcounter{utterance}  
    & & & \multicolumn{2}{p{0.3\linewidth}}{
        \cellcolor[rgb]{0.9,0.9,0.9}{
            \makecell[{{p{\linewidth}}}]{
                \texttt{\tiny{[GM$|$GM]}}
                \texttt{MOZGASD: P, (1,3)} \\
            }
        }
    }
    & & \\ \\

    \theutterance \stepcounter{utterance}  
    & & & \multicolumn{2}{p{0.3\linewidth}}{
        \cellcolor[rgb]{0.9,0.9,0.9}{
            \makecell[{{p{\linewidth}}}]{
                \texttt{\tiny{[GM$|$GM]}}
                \texttt{Büntetés: a (1,3) pozíció nem üres, tartalmazza: 'L'.} \\
\\ 
\texttt{Eddig összesen 7 büntetést kaptatok a megengedett 8{-}ból/ből.} \\
\texttt{Kérlek, próbáld újra!} \\
            }
        }
    }
    & & \\ \\

    \theutterance \stepcounter{utterance}  
    & & \multicolumn{4}{p{0.6\linewidth}}{
        \cellcolor[rgb]{0.9,0.9,0.9}{
            \makecell[{{p{\linewidth}}}]{
                \texttt{\tiny{[P1$\langle$GM]}}
                \texttt{Büntetés: a (1,3) pozíció nem üres, tartalmazza: 'L'.} \\
\\ 
\texttt{Eddig összesen 7 büntetést kaptatok a megengedett 8{-}ból/ből.} \\
\texttt{Kérlek, próbáld újra!} \\
            }
        }
    }
    & \\ \\

    \theutterance \stepcounter{utterance}  
    & \multicolumn{4}{p{0.6\linewidth}}{
        \cellcolor[rgb]{0.9,0.9,0.9}{
            \makecell[{{p{\linewidth}}}]{
                \texttt{\tiny{[P1$\rangle$GM]}}
                \texttt{MOZGASD: P, (1,4)} \\
            }
        }
    }
    & & \\ \\

    \theutterance \stepcounter{utterance}  
    & & & \multicolumn{2}{p{0.3\linewidth}}{
        \cellcolor[rgb]{0.9,0.9,0.9}{
            \makecell[{{p{\linewidth}}}]{
                \texttt{\tiny{[GM$|$GM]}}
                \texttt{MOZGASD: P, (1,4)} \\
            }
        }
    }
    & & \\ \\

    \theutterance \stepcounter{utterance}  
    & & & \multicolumn{2}{p{0.3\linewidth}}{
        \cellcolor[rgb]{0.9,0.9,0.9}{
            \makecell[{{p{\linewidth}}}]{
                \texttt{\tiny{[GM$|$GM]}}
                \texttt{Büntetés: a (1,4) pozíció nem üres, tartalmazza: '─'.} \\
\\ 
\texttt{Eddig összesen 8 büntetést kaptatok a megengedett 8{-}ból/ből.} \\
\texttt{Kérlek, próbáld újra!} \\
            }
        }
    }
    & & \\ \\

    \theutterance \stepcounter{utterance}  
    & & \multicolumn{4}{p{0.6\linewidth}}{
        \cellcolor[rgb]{0.9,0.9,0.9}{
            \makecell[{{p{\linewidth}}}]{
                \texttt{\tiny{[P1$\langle$GM]}}
                \texttt{Büntetés: a (1,4) pozíció nem üres, tartalmazza: '─'.} \\
\\ 
\texttt{Eddig összesen 8 büntetést kaptatok a megengedett 8{-}ból/ből.} \\
\texttt{Kérlek, próbáld újra!} \\
            }
        }
    }
    & \\ \\

    \theutterance \stepcounter{utterance}  
    & \multicolumn{4}{p{0.6\linewidth}}{
        \cellcolor[rgb]{0.9,0.9,0.9}{
            \makecell[{{p{\linewidth}}}]{
                \texttt{\tiny{[P1$\rangle$GM]}}
                \texttt{MONDD: Kész?} \\
            }
        }
    }
    & & \\ \\

    \theutterance \stepcounter{utterance}  
    & & & \multicolumn{2}{p{0.3\linewidth}}{
        \cellcolor[rgb]{0.9,0.9,0.9}{
            \makecell[{{p{\linewidth}}}]{
                \texttt{\tiny{[GM$|$GM]}}
                \texttt{MONDD: Kész?} \\
            }
        }
    }
    & & \\ \\

    \theutterance \stepcounter{utterance}  
    & & \multicolumn{4}{p{0.6\linewidth}}{
        \cellcolor[rgb]{0.9,0.9,0.9}{
            \makecell[{{p{\linewidth}}}]{
                \texttt{\tiny{[GM$\rangle$P2]}}
                \texttt{Az üzeneted továbbítva lett a másik játékosnak.} \\
\\ 
\\ 
\texttt{Jelenleg a(z) 7. körnél tartotok a maximum 12 körből.} \\
\texttt{Eddig összesen 8 büntetést kaptatok a megengedett 8{-}ból/ből.} \\
\texttt{A másik játékos ezt az üzenetet küldte:} \\
\\ 
\\ 
\texttt{"Kész?"} \\
\\ 
\\ 
\texttt{Mi legyen a következő parancsod?} \\
            }
        }
    }
    & \\ \\

    \theutterance \stepcounter{utterance}  
    & & & \multicolumn{4}{p{0.6\linewidth}}{
        \cellcolor[rgb]{0.9,0.9,0.9}{
            \makecell[{{p{\linewidth}}}]{
                \texttt{\tiny{[GM$\langle$P2]}}
                \texttt{MONDD: Nem, még nem vagyunk készen. Az L és C tárgyak pozíciója egyezik, de a P tárgy pozícióját még nem egyeztettük véglegesen. A mi P{-}nk az (1,3){-}on van. Kérlek, jelezd, hogy a te P{-}d is az (1,3){-}on van{-}e. Ha igen, akkor küldd a "kész!" üzenetet.} \\
            }
        }
    }
     \\ \\

    \theutterance \stepcounter{utterance}  
    & & & \multicolumn{2}{p{0.3\linewidth}}{
        \cellcolor[rgb]{0.9,0.9,0.9}{
            \makecell[{{p{\linewidth}}}]{
                \texttt{\tiny{[GM$|$GM]}}
                \texttt{MONDD: Nem, még nem vagyunk készen. Az L és C tárgyak pozíciója egyezik, de a P tárgy pozícióját még nem egyeztettük véglegesen. A mi P{-}nk az (1,3){-}on van. Kérlek, jelezd, hogy a te P{-}d is az (1,3){-}on van{-}e. Ha igen, akkor küldd a "kész!" üzenetet.} \\
            }
        }
    }
    & & \\ \\

    \theutterance \stepcounter{utterance}  
    & & \multicolumn{4}{p{0.6\linewidth}}{
        \cellcolor[rgb]{0.9,0.9,0.9}{
            \makecell[{{p{\linewidth}}}]{
                \texttt{\tiny{[P1$\langle$GM]}}
                \texttt{Az üzeneted továbbítva lett a másik játékosnak.} \\
\\ 
\\ 
\texttt{Jelenleg a(z) 7. körnél tartotok a maximum 12 körből.} \\
\texttt{Eddig összesen 8 büntetést kaptatok a megengedett 8{-}ból/ből.} \\
\texttt{A másik játékos ezt az üzenetet küldte:} \\
\\ 
\\ 
\texttt{"Nem, még nem vagyunk készen. Az L és C tárgyak pozíciója egyezik, de a P tárgy pozícióját még nem egyeztettük véglegesen. A mi P{-}nk az (1,3){-}on van. Kérlek, jelezd, hogy a te P{-}d is az (1,3){-}on van{-}e. Ha igen, akkor küldd a "kész!" üzenetet."} \\
\\ 
\\ 
\texttt{Mi legyen a következő parancsod?} \\
            }
        }
    }
    & \\ \\

    \theutterance \stepcounter{utterance}  
    & \multicolumn{4}{p{0.6\linewidth}}{
        \cellcolor[rgb]{0.9,0.9,0.9}{
            \makecell[{{p{\linewidth}}}]{
                \texttt{\tiny{[P1$\rangle$GM]}}
                \texttt{MONDD: Nem, a P nem az (1,3){-}on van, mert ott az L van. A P{-}t most az (1,4){-}re próbáltam mozgatni, de ott keret volt. Kérlek, mozgassd a saját P{-}det az (1,3){-}ról egy üres helyre, hogy szabaduljon az (1,3), vagy egyeztessünk egy másik közös pozíciót a P{-}nek. Például mozgassd a P{-}t az (1,4){-}re (ha nálad üres), és én is oda mozgatom.} \\
            }
        }
    }
    & & \\ \\

    \theutterance \stepcounter{utterance}  
    & & & \multicolumn{2}{p{0.3\linewidth}}{
        \cellcolor[rgb]{0.9,0.9,0.9}{
            \makecell[{{p{\linewidth}}}]{
                \texttt{\tiny{[GM$|$GM]}}
                \texttt{MONDD: Nem, a P nem az (1,3){-}on van, mert ott az L van. A P{-}t most az (1,4){-}re próbáltam mozgatni, de ott keret volt. Kérlek, mozgassd a saját P{-}det az (1,3){-}ról egy üres helyre, hogy szabaduljon az (1,3), vagy egyeztessünk egy másik közös pozíciót a P{-}nek. Például mozgassd a P{-}t az (1,4){-}re (ha nálad üres), és én is oda mozgatom.} \\
            }
        }
    }
    & & \\ \\

    \theutterance \stepcounter{utterance}  
    & & \multicolumn{4}{p{0.6\linewidth}}{
        \cellcolor[rgb]{0.9,0.9,0.9}{
            \makecell[{{p{\linewidth}}}]{
                \texttt{\tiny{[GM$\rangle$P2]}}
                \texttt{Az üzeneted továbbítva lett a másik játékosnak.} \\
\\ 
\\ 
\texttt{Jelenleg a(z) 8. körnél tartotok a maximum 12 körből.} \\
\texttt{Eddig összesen 8 büntetést kaptatok a megengedett 8{-}ból/ből.} \\
\texttt{A másik játékos ezt az üzenetet küldte:} \\
\\ 
\\ 
\texttt{"Nem, a P nem az (1,3){-}on van, mert ott az L van. A P{-}t most az (1,4){-}re próbáltam mozgatni, de ott keret volt. Kérlek, mozgassd a saját P{-}det az (1,3){-}ról egy üres helyre, hogy szabaduljon az (1,3), vagy egyeztessünk egy másik közös pozíciót a P{-}nek. Például mozgassd a P{-}t az (1,4){-}re (ha nálad üres), és én is oda mozgatom."} \\
\\ 
\\ 
\texttt{Mi legyen a következő parancsod?} \\
            }
        }
    }
    & \\ \\

    \theutterance \stepcounter{utterance}  
    & & & \multicolumn{4}{p{0.6\linewidth}}{
        \cellcolor[rgb]{0.9,0.9,0.9}{
            \makecell[{{p{\linewidth}}}]{
                \texttt{\tiny{[GM$\langle$P2]}}
                \texttt{MOZGASD: P, (1,4)} \\
            }
        }
    }
     \\ \\

    \theutterance \stepcounter{utterance}  
    & & & \multicolumn{2}{p{0.3\linewidth}}{
        \cellcolor[rgb]{0.9,0.9,0.9}{
            \makecell[{{p{\linewidth}}}]{
                \texttt{\tiny{[GM$|$GM]}}
                \texttt{MOZGASD: P, (1,4)} \\
            }
        }
    }
    & & \\ \\

    \theutterance \stepcounter{utterance}  
    & & & \multicolumn{2}{p{0.3\linewidth}}{
        \cellcolor[rgb]{0.9,0.9,0.9}{
            \makecell[{{p{\linewidth}}}]{
                \texttt{\tiny{[GM$|$GM]}}
                \texttt{Büntetés: a (1,4) pozíció nem üres, tartalmazza: '─'.} \\
\\ 
\texttt{Eddig összesen 9 büntetést kaptatok a megengedett 8{-}ból/ből.} \\
\texttt{Kérlek, próbáld újra!} \\
            }
        }
    }
    & & \\ \\

    \theutterance \stepcounter{utterance}  
    & & \multicolumn{4}{p{0.6\linewidth}}{
        \cellcolor[rgb]{0.9,0.9,0.9}{
            \makecell[{{p{\linewidth}}}]{
                \texttt{\tiny{[GM$\rangle$P2]}}
                \texttt{Büntetés: a (1,4) pozíció nem üres, tartalmazza: '─'.} \\
\\ 
\texttt{Eddig összesen 9 büntetést kaptatok a megengedett 8{-}ból/ből.} \\
\texttt{Kérlek, próbáld újra!} \\
            }
        }
    }
    & \\ \\

    \theutterance \stepcounter{utterance}  
    & & & \multicolumn{4}{p{0.6\linewidth}}{
        \cellcolor[rgb]{0.9,0.9,0.9}{
            \makecell[{{p{\linewidth}}}]{
                \texttt{\tiny{[GM$\langle$P2]}}
                \texttt{MONDD: A (1,4) nem üres, keret van ott. Nézzük meg a rácsot újra. Az L most az (1,1), a C az (5,1), a P az (1,3). Üres helyek: (1,2), (2,1), (3,1), (4,1), (6,1), (7,1), (2,2), (3,2), (4,2), (5,2), (6,2), (7,2), (2,3), (3,3), (4,3), (5,3), (6,3), (7,3), (2,4), (3,4), (4,4), (5,4), (6,4), (7,4), (2,5), (3,5), (4,5), (5,5), (6,5), (7,5), (2,6), (3,6), (4,6), (5,6), (6,6), (7,6), (2,7), (3,7), (4,7), (5,7), (6,7), (7,7). Javasolom, hogy a P{-}t mozgassuk az (2,2){-}re. Te is mozgassd a P{-}det az (2,2){-}re.} \\
            }
        }
    }
     \\ \\

    \theutterance \stepcounter{utterance}  
    & & & \multicolumn{2}{p{0.3\linewidth}}{
        \cellcolor[rgb]{0.9,0.9,0.9}{
            \makecell[{{p{\linewidth}}}]{
                \texttt{\tiny{[GM$|$GM]}}
                \texttt{MONDD: A (1,4) nem üres, keret van ott. Nézzük meg a rácsot újra. Az L most az (1,1), a C az (5,1), a P az (1,3). Üres helyek: (1,2), (2,1), (3,1), (4,1), (6,1), (7,1), (2,2), (3,2), (4,2), (5,2), (6,2), (7,2), (2,3), (3,3), (4,3), (5,3), (6,3), (7,3), (2,4), (3,4), (4,4), (5,4), (6,4), (7,4), (2,5), (3,5), (4,5), (5,5), (6,5), (7,5), (2,6), (3,6), (4,6), (5,6), (6,6), (7,6), (2,7), (3,7), (4,7), (5,7), (6,7), (7,7). Javasolom, hogy a P{-}t mozgassuk az (2,2){-}re. Te is mozgassd a P{-}det az (2,2){-}re.} \\
            }
        }
    }
    & & \\ \\

    \theutterance \stepcounter{utterance}  
    & & \multicolumn{4}{p{0.6\linewidth}}{
        \cellcolor[rgb]{0.9,0.9,0.9}{
            \makecell[{{p{\linewidth}}}]{
                \texttt{\tiny{[P1$\langle$GM]}}
                \texttt{Az üzeneted továbbítva lett a másik játékosnak.} \\
\\ 
\\ 
\texttt{Jelenleg a(z) 8. körnél tartotok a maximum 12 körből.} \\
\texttt{Eddig összesen 9 büntetést kaptatok a megengedett 8{-}ból/ből.} \\
\texttt{A másik játékos ezt az üzenetet küldte:} \\
\\ 
\\ 
\texttt{"A (1,4) nem üres, keret van ott. Nézzük meg a rácsot újra. Az L most az (1,1), a C az (5,1), a P az (1,3). Üres helyek: (1,2), (2,1), (3,1), (4,1), (6,1), (7,1), (2,2), (3,2), (4,2), (5,2), (6,2), (7,2), (2,3), (3,3), (4,3), (5,3), (6,3), (7,3), (2,4), (3,4), (4,4), (5,4), (6,4), (7,4), (2,5), (3,5), (4,5), (5,5), (6,5), (7,5), (2,6), (3,6), (4,6), (5,6), (6,6), (7,6), (2,7), (3,7), (4,7), (5,7), (6,7), (7,7). Javasolom, hogy a P{-}t mozgassuk az (2,2){-}re. Te is mozgassd a P{-}det az (2,2){-}re."} \\
\\ 
\\ 
\texttt{Mi legyen a következő parancsod?} \\
            }
        }
    }
    & \\ \\

    \theutterance \stepcounter{utterance}  
    & \multicolumn{4}{p{0.6\linewidth}}{
        \cellcolor[rgb]{0.9,0.9,0.9}{
            \makecell[{{p{\linewidth}}}]{
                \texttt{\tiny{[P1$\rangle$GM]}}
                \texttt{MOZGASD: P, (2,2)} \\
            }
        }
    }
    & & \\ \\

    \theutterance \stepcounter{utterance}  
    & & & \multicolumn{2}{p{0.3\linewidth}}{
        \cellcolor[rgb]{0.9,0.9,0.9}{
            \makecell[{{p{\linewidth}}}]{
                \texttt{\tiny{[GM$|$GM]}}
                \texttt{MOZGASD: P, (2,2)} \\
            }
        }
    }
    & & \\ \\

    \theutterance \stepcounter{utterance}  
    & & & \multicolumn{2}{p{0.3\linewidth}}{
        \cellcolor[rgb]{0.9,0.9,0.9}{
            \makecell[{{p{\linewidth}}}]{
                \texttt{\tiny{[GM$|$GM]}}
                \texttt{P moved to (2, 2)} \\
            }
        }
    }
    & & \\ \\

    \theutterance \stepcounter{utterance}  
    & & \multicolumn{4}{p{0.6\linewidth}}{
        \cellcolor[rgb]{0.9,0.9,0.9}{
            \makecell[{{p{\linewidth}}}]{
                \texttt{\tiny{[GM$\rangle$P2]}}
                \texttt{Az üzeneted továbbítva lett a másik játékosnak.} \\
\\ 
\\ 
\texttt{Jelenleg a(z) 9. körnél tartotok a maximum 12 körből.} \\
\texttt{Eddig összesen 9 büntetést kaptatok a megengedett 8{-}ból/ből.} \\
\texttt{A másik játékos mozgatott egy tárgyat a saját rácsán.} \\
\\ 
\\ 
\texttt{Mi legyen a következő parancsod?} \\
            }
        }
    }
    & \\ \\

    \theutterance \stepcounter{utterance}  
    & & & \multicolumn{4}{p{0.6\linewidth}}{
        \cellcolor[rgb]{0.9,0.9,0.9}{
            \makecell[{{p{\linewidth}}}]{
                \texttt{\tiny{[GM$\langle$P2]}}
                \texttt{MOZGASD: P, (2,2)} \\
            }
        }
    }
     \\ \\

    \theutterance \stepcounter{utterance}  
    & & & \multicolumn{2}{p{0.3\linewidth}}{
        \cellcolor[rgb]{0.9,0.9,0.9}{
            \makecell[{{p{\linewidth}}}]{
                \texttt{\tiny{[GM$|$GM]}}
                \texttt{MOZGASD: P, (2,2)} \\
            }
        }
    }
    & & \\ \\

    \theutterance \stepcounter{utterance}  
    & & & \multicolumn{2}{p{0.3\linewidth}}{
        \cellcolor[rgb]{0.9,0.9,0.9}{
            \makecell[{{p{\linewidth}}}]{
                \texttt{\tiny{[GM$|$GM]}}
                \texttt{P moved to (2, 2)} \\
            }
        }
    }
    & & \\ \\

    \theutterance \stepcounter{utterance}  
    & & \multicolumn{4}{p{0.6\linewidth}}{
        \cellcolor[rgb]{0.9,0.9,0.9}{
            \makecell[{{p{\linewidth}}}]{
                \texttt{\tiny{[P1$\langle$GM]}}
                \texttt{A(z) 'P' sikeresen át lett mozgatva ide: (2,2). A frissített rácsod így néz ki:} \\
\texttt{```} \\
\texttt{ 1234567} \\
\texttt{╔═══════╗} \\
\texttt{║◌◌◌◌C◌◌║ 1} \\
\texttt{║◌P◌◌◌◌◌║ 2} \\
\texttt{║L◌◌◌◌┌─╢ 3} \\
\texttt{╟──┐◌◌│◌║ 4} \\
\texttt{║◌◌├──┼─╢ 5} \\
\texttt{║◌◌│◌◌│◌║ 6} \\
\texttt{║◌◌│◌◌│◌║ 7} \\
\texttt{╚══╧══╧═╝} \\
\texttt{```} \\
\texttt{Jelenleg a(z) 9. körnél tartotok a maximum 12 körből.} \\
\texttt{Eddig összesen 9 büntetést kaptatok a megengedett 8{-}ból/ből.} \\
\texttt{A másik játékos mozgatott egy tárgyat a saját rácsán.} \\
\\ 
\\ 
\texttt{Mi legyen a következő parancsod?} \\
            }
        }
    }
    & \\ \\

    \theutterance \stepcounter{utterance}  
    & \multicolumn{4}{p{0.6\linewidth}}{
        \cellcolor[rgb]{0.9,0.9,0.9}{
            \makecell[{{p{\linewidth}}}]{
                \texttt{\tiny{[P1$\rangle$GM]}}
                \texttt{MONDD: Kész!} \\
            }
        }
    }
    & & \\ \\

    \theutterance \stepcounter{utterance}  
    & & & \multicolumn{2}{p{0.3\linewidth}}{
        \cellcolor[rgb]{0.9,0.9,0.9}{
            \makecell[{{p{\linewidth}}}]{
                \texttt{\tiny{[GM$|$GM]}}
                \texttt{MONDD: Kész!} \\
            }
        }
    }
    & & \\ \\

    \theutterance \stepcounter{utterance}  
    & & \multicolumn{4}{p{0.6\linewidth}}{
        \cellcolor[rgb]{0.9,0.9,0.9}{
            \makecell[{{p{\linewidth}}}]{
                \texttt{\tiny{[GM$\rangle$P2]}}
                \texttt{A(z) 'P' sikeresen át lett mozgatva ide: (2,2). A frissített rácsod így néz ki:} \\
\texttt{```} \\
\texttt{ 1234567} \\
\texttt{╔═══════╗} \\
\texttt{║L◌◌◌C◌◌║ 1} \\
\texttt{║◌P◌◌◌◌◌║ 2} \\
\texttt{║◌◌◌◌◌┌─╢ 3} \\
\texttt{╟──┐◌◌│◌║ 4} \\
\texttt{║◌◌├──┼─╢ 5} \\
\texttt{║◌◌│◌◌│◌║ 6} \\
\texttt{║◌◌│◌◌│◌║ 7} \\
\texttt{╚══╧══╧═╝} \\
\texttt{```} \\
\texttt{Jelenleg a(z) 10. körnél tartotok a maximum 12 körből.} \\
\texttt{Eddig összesen 9 büntetést kaptatok a megengedett 8{-}ból/ből.} \\
\texttt{A másik játékos ezt az üzenetet küldte:} \\
\\ 
\\ 
\texttt{"Kész!"} \\
\\ 
\\ 
\texttt{Mi legyen a következő parancsod?} \\
            }
        }
    }
    & \\ \\

    \theutterance \stepcounter{utterance}  
    & & & \multicolumn{4}{p{0.6\linewidth}}{
        \cellcolor[rgb]{0.9,0.9,0.9}{
            \makecell[{{p{\linewidth}}}]{
                \texttt{\tiny{[GM$\langle$P2]}}
                \texttt{MONDD: kész!} \\
            }
        }
    }
     \\ \\

    \theutterance \stepcounter{utterance}  
    & & & \multicolumn{2}{p{0.3\linewidth}}{
        \cellcolor[rgb]{0.9,0.9,0.9}{
            \makecell[{{p{\linewidth}}}]{
                \texttt{\tiny{[GM$|$GM]}}
                \texttt{MONDD: kész!} \\
            }
        }
    }
    & & \\ \\

    \theutterance \stepcounter{utterance}  
    & & \multicolumn{4}{p{0.6\linewidth}}{
        \cellcolor[rgb]{0.9,0.9,0.9}{
            \makecell[{{p{\linewidth}}}]{
                \texttt{\tiny{[P1$\langle$GM]}}
                \texttt{Az üzeneted továbbítva lett a másik játékosnak.} \\
\\ 
\\ 
\texttt{Jelenleg a(z) 10. körnél tartotok a maximum 12 körből.} \\
\texttt{Eddig összesen 9 büntetést kaptatok a megengedett 8{-}ból/ből.} \\
\texttt{A másik játékos ezt az üzenetet küldte:} \\
\\ 
\\ 
\texttt{"kész!"} \\
\\ 
\\ 
\texttt{Mi legyen a következő parancsod?} \\
            }
        }
    }
    & \\ \\

    \theutterance \stepcounter{utterance}  
    & \multicolumn{4}{p{0.6\linewidth}}{
        \cellcolor[rgb]{0.9,0.9,0.9}{
            \makecell[{{p{\linewidth}}}]{
                \texttt{\tiny{[P1$\rangle$GM]}}
                \texttt{A játék véget ért. Mindkét játékos megerősítette, hogy a tárgyak a megfelelő pozíciókon vannak (L: 1,1; C: 5,1; P: 2,2). A játékot megnyertétek, bár a büntetések száma elérte a maximumot.} \\
            }
        }
    }
    & & \\ \\

    \theutterance \stepcounter{utterance}  
    & & & \multicolumn{2}{p{0.3\linewidth}}{
        \cellcolor[rgb]{0.9,0.9,0.9}{
            \makecell[{{p{\linewidth}}}]{
                \texttt{\tiny{[GM$|$GM]}}
                \texttt{A játék véget ért. Mindkét játékos megerősítette, hogy a tárgyak a megfelelő pozíciókon vannak (L: 1,1; C: 5,1; P: 2,2). A játékot megnyertétek, bár a büntetések száma elérte a maximumot.} \\
            }
        }
    }
    & & \\ \\

    \theutterance \stepcounter{utterance}  
    & & & \multicolumn{2}{p{0.3\linewidth}}{
        \cellcolor[rgb]{0.9,0.9,0.9}{
            \makecell[{{p{\linewidth}}}]{
                \texttt{\tiny{[GM$|$GM]}}
                \texttt{Invalid response format} \\
            }
        }
    }
    & & \\ \\

    \theutterance \stepcounter{utterance}  
    & & & \multicolumn{2}{p{0.3\linewidth}}{
        \cellcolor[rgb]{0.9,0.9,0.9}{
            \makecell[{{p{\linewidth}}}]{
                \texttt{\tiny{[GM$|$GM]}}
                \texttt{Maximum number of penalties exceeded} \\
            }
        }
    }
    & & \\ \\

    \theutterance \stepcounter{utterance}  
    & & & \multicolumn{2}{p{0.3\linewidth}}{
        \cellcolor[rgb]{0.9,0.9,0.9}{
            \makecell[{{p{\linewidth}}}]{
                \texttt{\tiny{[GM$|$GM]}}
                \texttt{Initial states:} \\
\texttt{```} \\
\texttt{Player 1:} \\
\texttt{ 1234567} \\
\texttt{╔═══════╗} \\
\texttt{║◌◌◌◌L◌◌║ 1} \\
\texttt{║P◌◌◌◌◌◌║ 2} \\
\texttt{║◌C◌◌◌┌─╢ 3} \\
\texttt{╟──┐◌◌│◌║ 4} \\
\texttt{║◌◌├──┼─╢ 5} \\
\texttt{║◌◌│◌◌│◌║ 6} \\
\texttt{║◌◌│◌◌│◌║ 7} \\
\texttt{╚══╧══╧═╝} \\
\\ 
\\ 
\texttt{Player 2:} \\
\texttt{ 1234567} \\
\texttt{╔═══════╗} \\
\texttt{║◌◌◌◌C◌◌║ 1} \\
\texttt{║L◌◌◌◌◌◌║ 2} \\
\texttt{║P◌◌◌◌┌─╢ 3} \\
\texttt{╟──┐◌◌│◌║ 4} \\
\texttt{║◌◌├──┼─╢ 5} \\
\texttt{║◌◌│◌◌│◌║ 6} \\
\texttt{║◌◌│◌◌│◌║ 7} \\
\texttt{╚══╧══╧═╝} \\
\\ 
\texttt{```} \\
            }
        }
    }
    & & \\ \\

    \theutterance \stepcounter{utterance}  
    & & & \multicolumn{2}{p{0.3\linewidth}}{
        \cellcolor[rgb]{0.9,0.9,0.9}{
            \makecell[{{p{\linewidth}}}]{
                \texttt{\tiny{[GM$|$GM]}}
                \texttt{End states:} \\
\texttt{```} \\
\texttt{Player 1:} \\
\texttt{ 1234567} \\
\texttt{╔═══════╗} \\
\texttt{║◌◌◌◌C◌◌║ 1} \\
\texttt{║◌P◌◌◌◌◌║ 2} \\
\texttt{║L◌◌◌◌┌─╢ 3} \\
\texttt{╟──┐◌◌│◌║ 4} \\
\texttt{║◌◌├──┼─╢ 5} \\
\texttt{║◌◌│◌◌│◌║ 6} \\
\texttt{║◌◌│◌◌│◌║ 7} \\
\texttt{╚══╧══╧═╝} \\
\\ 
\\ 
\texttt{Player 2:} \\
\texttt{ 1234567} \\
\texttt{╔═══════╗} \\
\texttt{║L◌◌◌C◌◌║ 1} \\
\texttt{║◌P◌◌◌◌◌║ 2} \\
\texttt{║◌◌◌◌◌┌─╢ 3} \\
\texttt{╟──┐◌◌│◌║ 4} \\
\texttt{║◌◌├──┼─╢ 5} \\
\texttt{║◌◌│◌◌│◌║ 6} \\
\texttt{║◌◌│◌◌│◌║ 7} \\
\texttt{╚══╧══╧═╝} \\
\\ 
\texttt{```} \\
            }
        }
    }
    & & \\ \\

    \theutterance \stepcounter{utterance}  
    & & & \multicolumn{2}{p{0.3\linewidth}}{
        \cellcolor[rgb]{0.9,0.9,0.9}{
            \makecell[{{p{\linewidth}}}]{
                \texttt{\tiny{[GM$|$GM]}}
                \texttt{* success: False} \\
\texttt{* lose: False} \\
\texttt{* aborted: True} \\
\texttt{{-}{-}{-}{-}{-}{-}{-}} \\
\texttt{* Shifts: 3.00} \\
\texttt{* Max Shifts: 4.00} \\
\texttt{* Min Shifts: 2.00} \\
\texttt{* End Distance Sum: 2.00} \\
\texttt{* Init Distance Sum: 8.73} \\
\texttt{* Expected Distance Sum: 12.57} \\
\texttt{* Penalties: 10.00} \\
\texttt{* Max Penalties: 8.00} \\
\texttt{* Rounds: 10.00} \\
\texttt{* Max Rounds: 12.00} \\
\texttt{* Object Count: 3.00} \\
            }
        }
    }
    & & \\ \\

\end{supertabular}
}

\end{document}
